\chapter{Introduction}
\label{ch:introduction}

Digital signatures are a basic mechanism for public-key authentication.  A
signature scheme allows a signer with a secret key to authenticate messages so
that anyone with the corresponding public key can verify them.  The standard
adaptive security notion for signatures is existential unforgeability under
chosen-message attack, or \eufcma{}.  In this notion an adversary receives a
public key, adaptively queries a signing oracle, and wins if it outputs a valid
signature on a message that the signing oracle did not sign.

The post-quantum transition makes the assurance problem more difficult.  Modern
lattice-based signatures combine algebraic structure, randomized sampling,
Fiat-Shamir transformations, rejection sampling, and concrete hash functions.
Their proofs involve many games and many small probability losses.  A missing
freshness side condition or an incorrectly counted oracle query can invalidate a
security claim even when the high-level proof idea is correct.

This dissertation documents a machine-checked provable-security proof surface
for a formal \haetae{} random-oracle model in \easycrypt{}.  The qualifier
``formal model'' is essential: the checked surface is a conditional implication
about the EasyCrypt game model and its counted random-oracle interface, not a
separate machine-checked equivalence proof for every line of the February 2026
HAETAE specification or for a concrete SHAKE implementation.  It is
``machine checked'' because the games, modules, lemmas, and inequalities are
accepted by EasyCrypt rather than only argued in prose.  Its current
cryptographic reading must remain conditional: the final counted-ROM conclusion
depends on explicit hop and hardness premises, and one final hop premise is
vacuous under the present \ec{rejection\_sampling\_loss\_term} constants unless
that premise or loss normalization is strengthened in future proof work.

\section{Problem Statement}

The problem addressed by this work is the following:

\begin{quote}
How can the random-oracle-model proof of the formal \haetae{} scheme's
\eufcma{} security be represented, checked, maintained, and audited as an
EasyCrypt development whose provable-security surface is clearly separated from
unrelated generated certification artifacts, while the model/specification
boundary remains explicit?
\end{quote}

This formulation has two parts.  The first is mathematical: the proof must
connect the \eufcma{} game to the underlying assumptions through valid game
hops and reductions.  The second is engineering-oriented: the proof repository
must make it clear which EasyCrypt files are part of the security theorem and
how to re-run the proof.

\section{Contributions}

The dissertation records six contributions.

\begin{enumerate}[label=\textbf{C\arabic*.}]
\item A separated EasyCrypt proof-management folder for \haetae{} provable
security under \code{provable-security/}.
\item A manifest-driven verification gate that compiles only the files needed
for the formal security proof.
\item A structured account of the theorem chain from generic signature games to
\haetae{}-specific \eufcma{} security.
\item A detailed explanation of checked random-oracle programming, sampler
freshness, budgeted lifting, rejection sampling, and concrete bound algebra.
\item A dependency ledger separating checked theorem content from modeling
assumptions, audited specification correspondences, and out-of-scope
cryptanalytic or implementation claims.
\item A clear boundary between provable-security artifacts, modeled hash
interfaces, concrete FIPS202/SHAKE validation material, and generated
known-answer-test files.
\end{enumerate}

\section{Reading Guide}

\Cref{ch:background} reviews the cryptographic and formal background.
\Cref{ch:proof-scope} defines the checked proof surface.  Later chapters follow
the proof architecture: EasyCrypt framework, signature games, ROM programming,
sampler freshness, transcript hops, budgeted lifting, reductions, and
reproducibility.
